%! lualatex
\documentclass[a4paper,11pt]{ltjsarticle}

\title{ヨガポーズ採点 Web アプリ\\操作説明書}
\author{利用者向け（インストラクター・自習者）}
\date{2026-06-02 \quad 版 1.0}

% 共通プリアンブル（LuaLaTeX + 日本語）
% ビルド: lualatex 対象ファイル.tex（2回推奨）

\usepackage[hiragino-pron,deluxe]{luatexja-preset}
\usepackage{geometry}
\geometry{left=25mm,right=25mm,top=25mm,bottom=28mm}

\usepackage{graphicx}
\usepackage{booktabs}
\usepackage{longtable}
\usepackage{array}
\usepackage{hyperref}
\usepackage{xcolor}
\usepackage{listings}
\usepackage{fancyhdr}
\usepackage{enumitem}
\usepackage{tabularx}
\usepackage{fancyvrb}

\hypersetup{
  colorlinks=true,
  linkcolor=teal!60!black,
  urlcolor=teal!50!black,
  pdfborder={0 0 0},
}

% フォントは luatexja-preset（hiragino-pron）で設定済み
% Linux 等では common.tex 先頭を [noto-otf,deluxe] に変更してください
\setsansjfont{Hiragino Sans}
\setmonofont{Menlo}[Scale = 0.88]

\lstset{
  basicstyle = \ttfamily\small,
  breaklines = true,
  frame = single,
  rulecolor = \color{black!15},
  backgroundcolor = \color{black!3},
  columns = fullflexible,
  keepspaces = true,
  showstringspaces = false,
}

\newcommand{\doctypemark}{}
\newcommand{\setdoctype}[1]{\renewcommand{\doctypemark}{#1}}

\pagestyle{fancy}
\fancyhf{}
\fancyhead[L]{\small\leftmark}
\fancyhead[R]{\small\doctypemark}
\fancyfoot[C]{\thepage}
\renewcommand{\headrulewidth}{0.4pt}
\renewcommand{\footrulewidth}{0pt}

\setlist[itemize]{topsep=0.4em,itemsep=0.2em}
\setlist[enumerate]{topsep=0.4em,itemsep=0.2em}

\DefineVerbatimEnvironment{asciibox}{Verbatim}{%
  fontsize=\small,
  frame=single,
  rulecolor=\color{black!20},
  framesep=6pt,
  xleftmargin=1em,
  xrightmargin=1em,
}

\newcommand{\pkg}[1]{\texttt{#1}}
\newcommand{\code}[1]{\texttt{#1}}

\setdoctype{操作説明書}

\begin{document}
\maketitle
\thispagestyle{fancy}

\noindent
\textbf{本番 URL}: \url{https://yoga-pose-scorer.vercel.app}

\tableofcontents
\newpage

\section{このアプリでできること}

\begin{itemize}
  \item スマホや PC のカメラで、ヨガのポーズを\textbf{リアルタイム}で採点します
  \item \textbf{3 種類}のポーズに対応（木のポーズ / 戦士のポーズ2 / 下向きの犬）
  \item 映像は\textbf{あなたの端末の中だけ}で処理（サーバーに送りません）
\end{itemize}

採点はあくまで\textbf{目安}です。体調に合わせ、無理のない範囲でご利用ください。

\section{動作環境}

\begin{table}[h]
  \centering
  \begin{tabular}{@{}ll@{}}
    \toprule
    項目 & 要件 \\
    \midrule
    ブラウザ & Chrome、Safari、Edge、Firefox など最新版 \\
    接続 & インターネット（初回のみ AI モデルのダウンロード） \\
    カメラ & 内蔵または外付け Web カメラ \\
    HTTPS & 本番サイトは HTTPS のためカメラ利用可 \\
    \bottomrule
  \end{tabular}
\end{table}

\noindent\textbf{注意}: \code{http://} の一般サイトではカメラが使えないことがあります。本番 URL（\url{https://yoga-pose-scorer.vercel.app}）をご利用ください。

\section{画面の見方}

\begin{asciibox}
+---------------------------------------------+
|  ヨガポーズ採点                              |
|  （説明文）                                  |
+---------------------------------------------+
|  練習するポーズ  [v 木のポーズ    ]          |
|  ポーズの説明…                               |
+----------------------+----------------------+
|                      |   78 点              |
|   カメラ映像          |   良い               |
|   （骨格の線）        |   軸足… 85点         |
|                      |   上げた足… 72点     |
|  [カメラを開始]       |   （各項目の詳細）    |
+----------------------+----------------------+
\end{asciibox}

\begin{table}[h]
  \centering
  \begin{tabular}{@{}ll@{}}
    \toprule
    場所 & 内容 \\
    \midrule
    上 & タイトル・ポーズの選び方 \\
    左（大きい枠） & カメラ映像。骨格が緑色の線で重なる \\
    右のパネル & \textbf{総合スコア}と\textbf{項目ごとの点数・フィードバック} \\
    下 & 利用上の注意 \\
    \bottomrule
  \end{tabular}
\end{table}

スマホなど狭い画面では、\textbf{右のパネルが下に表示}されます。スコアが見えないときは\textbf{下にスクロール}してください。

\section{基本的な使い方}

\subsection{ステップ 1 --- サイトを開く}

ブラウザで次の URL を開きます。

\noindent\textbf{\url{https://yoga-pose-scorer.vercel.app}}

\subsection{ステップ 2 --- ポーズを選ぶ}

「\textbf{練習するポーズ}」のプルダウンから練習したいポーズを選びます。

\begin{table}[h]
  \centering
  \begin{tabular}{@{}ll@{}}
    \toprule
    ポーズ名 & ざっくりした形 \\
    \midrule
    木のポーズ & 片足で立ち、もう一方の足を内ももに \\
    戦士のポーズ2 & 前足の膝を曲げ、腕を左右に伸ばす \\
    下向きの犬 & 手と足で床を押し、お尻を上げた逆 V 字 \\
    \bottomrule
  \end{tabular}
\end{table}

選ぶと、その下に短い説明が表示されます。

\subsection{ステップ 3 --- カメラを開始する}

\begin{enumerate}
  \item \textbf{「カメラを開始」}をタップ／クリック
  \item ブラウザから\textbf{カメラの使用を許可}
  \item 初回のみ\textbf{「モデル読み込み中…」}と表示されることがある（数十秒の場合あり）
\end{enumerate}

許可を拒否した場合は、ブラウザの設定からこのサイトのカメラを許可し、ページを再読み込みしてください。

\subsection{ステップ 4 --- ポーズをとる}

\begin{itemize}
  \item カメラから\textbf{2--3\,m 程度}離れ、\textbf{頭から足まで}が画面に入るようにする
  \item 明るい場所で、体のラインが隠れない服装がおすすめ
  \item 画面は鏡のように左右反転（見慣れた向き）
\end{itemize}

骨格の線（緑）が体に沿って表示されれば、検出できています。

\subsection{ステップ 5 --- スコアを確認する}

\textbf{右側のパネル}（スマホでは下）を見ます。

\begin{table}[h]
  \centering
  \begin{tabular}{@{}ll@{}}
    \toprule
    表示 & 意味 \\
    \midrule
    大きい数字 + 「点」 & 総合スコア（0--100） \\
    その下の短い文 & 総合評価（例: 素晴らしい / 良い） \\
    リストの各項目 & 部位ごとの点数 \\
    「現在 ○° / 目標 ○° --- …」 & 今の角度とアドバイス \\
    \bottomrule
  \end{tabular}
\end{table}

総合スコアが\textbf{「---」}のまま、「全身を映してください」と出る場合は、距離やカメラの向きを調整してください。

\subsection{ステップ 6 --- 終了する}

\textbf{「カメラを停止」}でカメラをオフにできます。ブラウザのタブを閉じてもカメラは停止します。

\section{スコア・フィードバックの読み方}

\subsection{総合スコアの目安}

\begin{table}[h]
  \centering
  \begin{tabular}{@{}ll@{}}
    \toprule
    点数 & 表示される評価 \\
    \midrule
    85 点以上 & 素晴らしい \\
    70 点以上 & 良い \\
    50 点以上 & もう一息 \\
    50 点未満 & 調整しましょう \\
    \bottomrule
  \end{tabular}
\end{table}

\subsection{項目ごとの表示例}

\begin{asciibox}
軸足（左）の伸展          85点
現在 172° / 目標 175° — とても良いです
\end{asciibox}

\begin{table}[h]
  \centering
  \begin{tabular}{@{}ll@{}}
    \toprule
    部分 & 説明 \\
    \midrule
    現在 ○° & カメラが検出した今の角度 \\
    目標 ○° & このポーズで近づけたい角度 \\
    最後の文 & 調整のヒント \\
    \bottomrule
  \end{tabular}
\end{table}

よく出るヒント:

\begin{table}[h]
  \centering
  \begin{tabular}{@{}ll@{}}
    \toprule
    ヒント & 意味 \\
    \midrule
    とても良いです & 目標にかなり近い \\
    もう少し調整してみてください & あと少し \\
    角度を少し大きく / 小さく & 関節をもう少し曲げる／伸ばす \\
    \bottomrule
  \end{tabular}
\end{table}

\subsection{左右について}

現在のバージョンでは、\textbf{左足・右足がポーズごとに固定}されています（例: 木のポーズは「左が軸足」想定）。反対の足で行うと、項目名と実際の足が一致しないことがありますが、角度の数値は参考にできます。

\section{うまくいかないとき}

\begin{table}[h]
  \centering
  \small
  \begin{tabular}{@{}p{0.32\linewidth}p{0.58\linewidth}@{}}
    \toprule
    症状 & 対処 \\
    \midrule
    カメラが真っ暗 & 許可ダイアログで「許可」／設定 $\to$ プライバシー $\to$ カメラ \\
    ずっと「モデル読み込み中」 & 回線を確認し、再読み込み（F5） \\
    スコアが「---」のまま & 全身が入るよう距離・明るさを調整 \\
    骨格線が乱れる & 逆光を避ける、体に物が重ならないようにする \\
    点数が急に上下する & 正常。姿勢をキープすると安定 \\
    初回だけ遅い & AI モデルのダウンロード。2 回目以降は速い \\
    \bottomrule
  \end{tabular}
\end{table}

\section{プライバシーについて}

\begin{itemize}
  \item カメラ映像は\textbf{あなたのブラウザ内}で処理
  \item 当アプリのサーバーに\textbf{映像や録画を保存しない}
  \item 初回に Google / jsDelivr 経由で AI モデルファイルを取得
\end{itemize}

\section{開発者向け（ローカルで試す場合）}

\begin{lstlisting}
cd /Users/waka/yoga-pose-scorer
npm install
npm run dev
\end{lstlisting}

ブラウザで \url{http://localhost:5173} を開きます。技術詳細は仕様書を参照してください。

\section{お問い合わせ・フィードバック}

仕様変更やポーズ追加の要望は、プロジェクト管理者（リポジトリオーナー）までご連絡ください。

\section{改訂履歴}

\begin{table}[h]
  \centering
  \begin{tabular}{@{}lll@{}}
    \toprule
    日付 & 版 & 内容 \\
    \midrule
    2026-06-02 & 1.0 & 初版 \\
    \bottomrule
  \end{tabular}
\end{table}

\end{document}
